\chapter{Fundamentação Teórica}
\label{cap:fundamentacao}

% -----------------------------------------------------------------
% Fio argumentativo do capítulo (funil até a lacuna):
%
%   2.1  Estabelece que o problema é real e multifatorial.
%   2.2  Mostra a tradição de apoio à decisão técnica e suas
%        limitações diante de dados sensoriais modernos.
%   2.3  Apresenta a primeira família de técnicas (anomalias
%        em séries operacionais) que sustenta o Pilar EQ. RQ1.
%   2.4  Apresenta a segunda família (estatística histórica por
%        lote, fornecedor, máquina) que sustenta o Pilar MP. RQ3.
%   2.5  Discute a integração das duas famílias e a
%        explicabilidade, apontando a lacuna que o WireSense
%        ocupa. RQ2 e RQ4.
%
% Protocolo da revisão sistemática:
%   docs/dissertacao/protocolo_revisao.md
% Sementes para bootstrap:
%   Chandola et al. (2009), Pang et al. (2021), Liu et al. (2008).
% -----------------------------------------------------------------


\section{Rompimentos em processos industriais}
\label{sec:rompimentos}

% Função retórica: estabelecer que o problema é real e
% multifatorial.
% Conteúdo a desenvolver:
%   - Natureza multifatorial do evento (matéria-prima,
%     equipamento, regulagem, interação).
%   - Efeitos em produtividade, qualidade e custo.
%   - Específico ao domínio: defeitos de vergalhão, falhas de
%     trefilação, esmaltação e bobinagem. RQ5.
% Conclui apontando que abordagens monofatoriais são
% insuficientes para diagnosticar a origem provável.


\section{Diagnóstico de falhas, manutenção e qualidade}
\label{sec:diagnostico}

% Função retórica: situar o trabalho na tradição de apoio à
% decisão técnica em manufatura.
% Conteúdo a desenvolver:
%   - Análise de falhas em ambiente produtivo (FMEA, RCA,
%     árvore de falhas, controle estatístico de processo).
%   - Limitações dessas ferramentas para absorver dado
%     sensorial em alta frequência.
% Conclui apontando a necessidade de incorporar dados
% operacionais e históricos de forma sistemática.


\section{Detecção de anomalias em dados industriais}
\label{sec:anomalias}

% Função retórica: apresentar a família de técnicas que
% sustenta o Pilar EQ.
% RQ1.
% Conteúdo a desenvolver:
%   - Conceito de anomalia em séries temporais multivariadas.
%   - Aprendizado não supervisionado, com destaque para
%     Isolation Forest e variantes.
%   - Métricas de avaliação em ausência de rótulo.
%   - Aplicações em manufatura e limites conhecidos.
% Conclui apontando que ML isolado não distingue causa
% associada a matéria-prima e que falta integração com
% conhecimento histórico de lote e fornecedor.


\section{Análise estatística histórica de falhas}
\label{sec:estatistica_historica}

% Função retórica: apresentar a família de técnicas que
% sustenta o Pilar MP.
% RQ3.
% Conteúdo a desenvolver:
%   - Taxa de falha, exposição operacional e risco relativo.
%   - Recorrência por lote, material, fornecedor e máquina.
%   - Efeito de lote e variabilidade entre fornecedores.
%   - Indicadores de qualidade de matéria-prima em
%     processo contínuo.
% Conclui apontando que estatística histórica isolada não
% captura desvio operacional em janela curta e depende de
% agregação que pode atrasar o diagnóstico.


\section{Sistemas híbridos e explicabilidade}
\label{sec:hibridos_xai}

% Função retórica: mostrar que combinar as duas famílias é a
% direção promissora, mas pouco consolidada no domínio de fios
% esmaltados, e apontar a lacuna.
% RQ2 e RQ4.
% Conteúdo a desenvolver:
%   - Estratégias de combinação (regras, ensemble, fusão de
%     evidências, sistemas em camadas).
%   - Importância de resultados interpretáveis em chão de
%     fábrica.
%   - Interpretabilidade aplicada a manutenção, qualidade e
%     engenharia.
% Conclui com o parágrafo da lacuna:
%   - Existe maturidade nas duas famílias isoladas.
%   - Há trabalhos pontuais de integração em outros domínios.
%   - Falta uma arquitetura híbrida e explicável aplicada ao
%     domínio específico de fios esmaltados, com regra de
%     cruzamento entre máquinas para diferenciar falha de
%     equipamento de suspeita de matéria-prima.
%   - Essa lacuna motiva o capítulo seguinte.
